\documentclass{article}
\usepackage{amsmath}
\usepackage{amssymb}
\usepackage{algorithm}
\usepackage{algpseudocode}
\usepackage{geometry}
 \geometry{
 a4paper,
 total={170mm,257mm},
 left=20mm,
 top=20mm,
 }

\title{Multi-Line Survival Data Generation Process}
\author{}
\date{}

\begin{document}
This document describes the data generation process for a multi-line survival model that simulates sequential treatment episodes for patients, incorporating autoregressive feature dynamics, disease progression, cumulative toxicity effects, and time-varying censoring.
Let's denote by: 
\begin{itemize}
    \item $X_{i,\ell} = (X^{j}_{i,\ell})_{j \leq J}$ the feature vector of patient $i$ at treatment line l 
    \item $T_{i,\ell}$ the treatment assigned at line $l$
    \item $Y_{i,\ell}$ the observed time (either event or censoring) at line $l$
    \item $E_{i,\ell}$ the event indicator. 
\end{itemize}
Features evolve across treatment lines following an AR($p$)
\begin{align*}
    &X_{i,0} \sim \mathcal{N}(\mu, 1), \quad j \in \{1, \ldots, J\}\\
    &X_{i,\ell} = \sum_{k=1}^{p} \alpha_k X_{i,\ell-k} + \epsilon_\ell, \quad \epsilon_\ell \sim \mathcal{N}(0, \sigma I_p)
\end{align*}
At each line $l$, treatment is randomly assigned among $r$ options with probabilities $\{p_1, \ldots, p_{r}\}$:
\begin{equation*}
    T_{i,\ell} \sim \mathcal{M}(1, \{p_1, p_2, \ldots, p_r\})
\end{equation*}
For each treatment, we attribute specific survival and toxicity parameters:
\begin{itemize}
    \item Survival parameters: $(k_T, \lambda_T)$ (Weibull shape and scale)
    \item Toxicity weight: $w_T \in [0, 1]$
\end{itemize}
Toxicity accumulates across previous treatment lines:
\begin{equation*}
    \text{Tox}_{i,\ell} \mid \{Y_{i,j}, T_{i,j}\}_{j < \ell} = \sum_{j=1}^{\ell-1} w_{T_{i,j}} \cdot Y_{i,j}
\end{equation*}
We can therefore define the hazard function at line $l$ for patient $i$ as:
\begin{equation}
    \eta_{i,\ell} = X_{i,\ell}^\top \beta + \beta_{\text{line}} \cdot \ell + \beta_{\text{tox}} \cdot \text{Tox}_{i,\ell}
\end{equation}
where:
\begin{itemize}
    \item $\beta \in \mathbb{R}^p$: feature coefficients
    \item $\beta_{\text{line}}$: disease progression coefficient
    \item $\beta_{\text{tox}}$: cumulative toxicity coefficient
\end{itemize}
In order to model treatment burden and disease progression, we define a baseline hazard that increases with line number:
$$
    h_{\text{base}}(\ell) = \exp(\rho \cdot \ell)
$$
where $\rho$ is the progression rate (increased hazard per line).
\paragraph{Event time generation:}
The event time follows a Weibull distribution with treatment-specific parameters and covariate adjustments:
$$ 
    Y_{i,\ell} \mid (T_{i,\ell}, X_{i,\ell}, Y_{i,\ell})_{i<l} = \frac{\exp(-\eta_{i,\ell})}{h_{\text{base}}(\ell)} \cdot W_{i,\ell}
$$
where $W_\ell \sim \text{Weibull}(k_{T_{i,\ell}}, \lambda_{T_{i,\ell}})$
\paragraph{Censoring time generation:}
Censoring time increases with line number:
\begin{align*}
    &\lambda_{\text{cens}}(\ell) = \lambda_0 (1 + \gamma \cdot \ell) \\
    &T_{i,\ell}^{\text{cens}} \sim \text{Exponential}(\lambda_{\text{cens}}(\ell))
\end{align*}    
where $\lambda_0$ is the base censoring rate and $\gamma$ is the per-line increase rate.
\paragraph{Observed data:}
The observed time and event indicator are:
\begin{align*}
    Y_{i,\ell} &= \min(T_{i,\ell}^{\text{event}}, T_{i,\ell}^{\text{cens}}) \\
    E_{i,\ell} &= \mathbb{I}(T_{i,\ell}^{\text{event}} \leq T_{i,\ell}^{\text{cens}})
\end{align*}
\begin{algorithm}
\caption{Multi-Line Survival Data Generation}
\begin{algorithmic}[1]
\For{each patient $i = 1, \ldots, n$}
    \State Sample $L_i \sim \text{Uniform}\{1, \ldots, L_{\max}\}$
    \State Initialize $X_{i,0}, \ldots, X_{i,p-1} \sim \mathcal{N}(0, I_p)$
    \State Generate AR($p$) process for $B + L_i$ steps
    \State Discard first $B$ burn-in steps
    \State Initialize: $\text{Tox}_{i,1} = 0$, previous treatments = []
    \For{each line $\ell = 1, \ldots, L_i$}
        \State Sample treatment $T_{i,\ell} \sim \text{Categorical}(\{A, B, C\})$
        \State Compute $\eta_{i,\ell}$ using Eq. (6)
        \State Compute $h_{\text{base}}(\ell)$ using Eq. (7)
        \State Generate $T_{i,\ell}^{\text{event}}$ using Eq. (8)
        \State Generate $T_{i,\ell}^{\text{cens}}$ using Eq. (10)
        \State Set $Y_{i,\ell} = \min(T_{i,\ell}^{\text{event}}, T_{i,\ell}^{\text{cens}})$
        \State Set $E_{i,\ell} = \mathbb{I}(T_{i,\ell}^{\text{event}} \leq T_{i,\ell}^{\text{cens}})$
        \State Update cumulative toxicity: $\text{Tox}_{i,\ell+1} = \text{Tox}_{i,\ell} + w_{T_{i,\ell}} \cdot Y_{i,\ell}$
    \EndFor
\EndFor
\end{algorithmic}
\end{algorithm}






\section{Bias and Convergence Properties of the CoxTime Gradient Estimator}

\subsection{Setup}

Consider the Cox partial log-likelihood
\[
\ell(\theta)
= 
\sum_{i: D_i = 1}
\left[
g_\theta(x_i)
-
\log \Big( \sum_{j \in R_i} \exp(g_\theta(x_j)) \Big)
\right],
\]
where $R_i$ is the risk set for individual $i$.

The gradient of the negative log-likelihood is
\[
\nabla L(\theta)
=
\sum_{i : D_i = 1}
\left(
\nabla g_\theta(x_i)
-
\frac{\sum_{j \in R_i} \exp(g_\theta(x_j)) \nabla g_\theta(x_j)}
{\sum_{j \in R_i} \exp(g_\theta(x_j))}
\right).
\]

This denominator requires summing over the \textit{full risk set}, which prevents minibatching.

CoxTime replaces $R_i$ with a random subset $\widetilde{R}_i$ of fixed size $m$, typically $m = |\text{batch}|$:
\[
\widetilde{R}_i \subseteq R_i,\quad |\widetilde{R}_i| = m,\quad i\in \widetilde{R}_i.
\]

The approximate loss is
\[
\widetilde{L}(\theta)
=
\sum_{i : D_i = 1}
\log \left(
\sum_{j \in \widetilde{R}_i}
\exp(g_\theta(x_j) - g_\theta(x_i))
\right).
\]

\subsection{Bias in the CoxTime Gradient}

The true gradient term for event $i$ is
\[
G_i(\theta)
=
\mathbb{E}_{J \sim p_i}
\big[
\nabla g_\theta(x_J)
\big],
\qquad
p_i(j) = \frac{\exp(g_\theta(x_j))}{\sum_{k \in R_i} \exp(g_\theta(x_k))}.
\]

The CoxTime estimator replaces $R_i$ by a subsample $\widetilde{R}_i$ with uniform sampling.  
Thus, the estimator uses
\[
\widetilde{G}_i(\theta)
=
\frac{\sum_{j \in \widetilde{R}_i}
\exp(g_\theta(x_j))\,\nabla g_\theta(x_j)}
{\sum_{j \in \widetilde{R}_i}
\exp(g_\theta(x_j))}.
\]

Since the sampling distribution is uniform rather than proportional to $\exp(g)$, we have:

\[
\mathbb{E}[\widetilde{G}_i(\theta)]
\neq
G_i(\theta)
\quad\text{(in general)}.
\]

The bias is:
\[
\mathrm{Bias}(\widetilde{G}_i) 
=
\mathbb{E}[\widetilde{G}_i] - G_i.
\]

\subsection{Magnitude of the Bias}

Let $Z_{i} = \sum_{j \in R_i} \exp(g_\theta(x_j))$ and  
$Z_{i,m} = \sum_{j \in \widetilde{R}_i} \exp(g_\theta(x_j))$.

If the subsampling is uniform without replacement:
\[
\mathbb{E}[Z_{i,m}] 
= \frac{m}{|R_i|} Z_i.
\]

Using Taylor approximations of the ratio of random variables:
\[
\frac{A}{B} = 
\frac{\mathbb{E}[A]}{\mathbb{E}[B]}
\left(
1 + O\!\left(\frac{\mathrm{Var}(A)}{A^2}\right)
+
O\!\left(\frac{\mathrm{Var}(B)}{B^2}\right)
\right).
\]

Thus the bias is of order
\[
\|\mathrm{Bias}(\widetilde{G}_i)\|
=
O\left(
\frac{1}{m}
\right),
\]
assuming $\exp(g_\theta(x_j))$ has bounded variance.

Therefore:
- Small batches $\Rightarrow$ significant bias  
- Large batches $\Rightarrow$ bias shrinks as $1/m$  
- Full batch $\Rightarrow$ zero bias  

\subsection{Unbiased or Less Biased Alternatives}

\paragraph{1. Weighted subsampling (Horvitz–Thompson).}

If every sampled $j$ is reweighted by the inverse selection probability:
\[
w_j = \frac{1}{p(j)} = \frac{|R_i|}{m},
\]
then the estimator
\[
\widehat{Z}_i = \sum_{j \in \widetilde{R}_i} w_j \exp(g_\theta(x_j))
\]
is unbiased:
\[
\mathbb{E}[\widehat{Z}_i] = Z_i.
\]

Same for the numerator.  
Therefore the ratio has bias only of order
\[
O\!\left(\frac{1}{m}\right)
\quad\text{instead of the naive } O(1).
\]

\paragraph{2. Importance sampling.}

Alternatively, sample $j$ with probability proportional to $\exp(g(x_j))$.

Then the gradient estimator is *exactly unbiased*:
\[
\mathbb{E}[\widetilde{G}_i]=G_i.
\]

The difficulty is that this sampling requires knowing all $\exp(g(x_j))$, which again breaks minibatching.

\subsection{Convergence of SGD with Biased Gradients}

Standard theory (Bottou, 2018) guarantees convergence if:

\begin{itemize}
\item The bias is bounded:
\[
\|\mathbb{E}[\widetilde{G}(\theta)] - \nabla L(\theta)\|
\le B.
\]
\item The learning rate decays: 
\[
\sum_t \eta_t = \infty,\qquad \sum_t \eta_t^2 < \infty.
\]
\item The objective is smooth.
\end{itemize}

The SGD iterates converge to a neighborhood of the optimum of size:
\[
O(B).
\]

Thus:
- DeepSurv (full risk set $B=0$) converges to the true optimum.  
- CoxTime (subset $B>0$) converges near the optimum, and the error shrinks as $O(1/m)$.

\subsection{Summary}

\begin{itemize}
\item CoxTime uses minibatch subsampling of the risk set, making the gradient estimator biased.
\item The bias is of order $1/m$ and vanishes as the sampled risk set grows.
\item Weighted sampling (Horvitz--Thompson) reduces the bias.
\item Importance sampling can make the estimator unbiased but is computationally infeasible.
\item SGD with small bias converges to a neighborhood of the optimum.
\end{itemize}

\end{document}